\section{Problem Summary}
An interval \([l,r]\) is transformed into
\[
f([l,r]) = [\min(a_l,\dots,a_r), \max(a_l,\dots,a_r)].
\]
For each query interval, we need the minimum number of repeated applications of \(f\) needed to reach \([1,n]\), or report that this never happens.

\section{Editorial}
\subsection{Motivation}
The function is not local in the usual sense: one application can move the interval far away, and naive simulation for every query is far too slow.

The hard part is finding a state description that can be lifted by powers of two.

\subsection{Key Observations}
\begin{itemize}
\item For a singleton interval \([i,i]\), the next state is just \([a_i,a_i]\). So singleton states are simple function composition on indices.
\item The nontrivial states are intervals of length at least \(2\). For those, it is enough to understand what happens to adjacent pairs \([i,i+1]\).
\item If two source intervals overlap at some index \(x\), then after applying \(f\), both resulting intervals contain \(a_x\). So overlapping source intervals always map to overlapping destination intervals.
\item Because of that overlap property, the images of consecutive adjacent pairs remain a chain of overlapping intervals at every power of two.
\end{itemize}

\subsection{Derivation}
Define
\[
g_t(i) = f^{2^t}([i,i+1])
\]
for \(1 \le i < n\), and
\[
p_t(i) = f^{2^t}([i,i])
\]
for singleton states.

The singleton transition is immediate:
\[
p_0(i) = [a_i,a_i], \qquad p_{t+1}(i) = p_t(a_i) \text{ on the index level}.
\]
So we can store only the index reached by a singleton after \(2^t\) steps.

Now look at any interval \([l,r]\) with \(l < r\). After one step,
\[
f([l,r]) = \bigcup_{i=l}^{r-1} f([i,i+1])
\]
in the interval sense: the minimum over the whole segment is the minimum of the pair-minima, and similarly for the maximum.

The same remains true after every power of two. Inductively, the intervals \(g_t(l), g_t(l+1), \dots, g_t(r-1)\) overlap consecutively, so their union is again one interval. Applying \(f\) to that union is the same as taking the hull of the next-step images of those pair intervals. Therefore
\[
f^{2^t}([l,r]) = \operatorname{hull}_{i=l}^{r-1} g_t(i).
\]

That identity is the whole problem. Once every \(g_t(i)\) is known, any jump by \(2^t\) steps becomes:
\begin{itemize}
\item a singleton jump if the current interval has length \(1\)
\item one range-hull query over the array \(g_t\) otherwise
\end{itemize}

\subsection{Algorithm}
\begin{enumerate}
\item Precompute singleton jumps \(p_t\) by repeated composition.
\item Precompute pair jumps \(g_t\).
  To build \(g_{t+1}(i)\):
  \begin{itemize}
  \item if \(g_t(i)\) is already a singleton \([x,x]\), then \(g_{t+1}(i) = p_t(x)\)
  \item otherwise \(g_{t+1}(i)\) is the hull of \(g_t(x)\) for all adjacent pairs inside that interval, so one range query on level \(t\) is enough
  \end{itemize}
\item Use a sparse table on each level to answer those hull queries in \(O(1)\).
\item For each query, binary-lift on the number of applications:
  keep the current interval, try to add powers of two from large to small, and only accept a jump if it still does \emph{not} reach \([1,n]\).
\end{enumerate}

\subsection{Why It Works}
The preprocessing is correct because it stores exactly the result of \(2^t\) applications for every adjacent pair and every singleton.

The range-hull formula is correct because consecutive pair-images always overlap, so their union stays an interval, and applying \(f\) to that interval only depends on the minimum left endpoint and the maximum right endpoint among the next-step pieces.

Binary lifting works because once \([1,n]\) is reached, it becomes a fixed point. Reaching \([1,n]\) means the array contains both values \(1\) and \(n\), hence the minimum and maximum over the full interval are still \(1\) and \(n\). So the predicate ``already reached \([1,n]\)'' is monotone in the number of steps.

\section{Pseudocode}
Read \(n\), the array, and the queries.

Precompute singleton jumps for powers of two.
Precompute pair jumps for powers of two:
\begin{itemize}
\item build a sparse table for the current pair level
\item for each adjacent pair image:
  \begin{itemize}
  \item if it is a singleton, advance it with the singleton table
  \item otherwise query the hull of all pair images inside that interval
  \end{itemize}
\end{itemize}

For each query:
\begin{itemize}
\item if it is already \([1,n]\), print \(0\)
\item if it is a singleton and \(n > 1\), print \(-1\)
\item otherwise try powers of two from large to small
\item keep every jump that still avoids \([1,n]\)
\item after the loop, check whether one more step reaches \([1,n]\)
\end{itemize}

\section{Complexity Analysis}
\begin{itemize}
\item Preprocessing pair levels with sparse tables: \(O(n \log^2 n)\)
\item Query phase: \(O(q \log n)\) once the levels are ready, plus another \(O(n \log^2 n)\) for the transient sparse tables used during binary lifting
\item Memory: \(O(n \log n)\)
\end{itemize}

\section{Notes}
\begin{itemize}
\item A singleton query can never become \([1,n]\) when \(n > 1\), because singleton states always stay singletons.
\item The implementation uses \(34\) lifting levels. That is safely above the number of distinct interval states for \(n \le 10^5\).
\item The merge operation is just interval hull: take the minimum left endpoint and the maximum right endpoint.
\end{itemize}
